% ============ 9 上同调: 定义与基本性质 (原书 pp.64-69) ============
\chapter{上同调：定义与基本性质}\label{sec:9}
\index{上同调}
\index{导出函子}
\index{Ext 群}
\index{带支集上同调}
\index{切除}
\index{同伦公理}
\index{有理等价}
\index{维数公理}

在上一节中，我们证明了 Abel 群层范畴 $\mathrm{Sh}(X_{\mathrm{et}})$ 是具有足够内射对象的 Abel 范畴.函子
\[
F \mapsto \Gamma(X, F) \colon \mathrm{Sh}(X_{\mathrm{et}}) \to \mathbf{Ab}
\]
左正合，我们定义 $H^{r}(X_{\mathrm{et}}, \Lambda)$ 为它的第 $r$ 个右导出函子.显式地，对层 $F$，选取内射分解
\[
0 \to F \to \mathcal{I}^{0} \to \mathcal{I}^{1} \to \mathcal{I}^{2} \to \cdots,
\]
并对它应用函子 $\Gamma(X, \cdot)$ 得到复形
\[
\Gamma(X, \mathcal{I}^{0}) \to \Gamma(X, \mathcal{I}^{1}) \to \Gamma(X, \mathcal{I}^{2}) \to \cdots.
\]
它(一般地)不再正合，而 $H^{r}(X_{\mathrm{et}}, F)$ 定义为它的第 $r$ 个上同调群.导出函子理论表明：
\begin{enumerate}
  \item 对任意层 $F$，$H^{0}(X_{\mathrm{et}}, F) = \Gamma(X, F)$；
  \item 若 $\mathcal{I}$ 内射，则对 $r > 0$，$H^{r}(X_{\mathrm{et}}, \mathcal{I}) = 0$；
  \item 层的短正合序列
  \[
  0 \to F' \to F \to F'' \to 0
  \]
  产生长正合序列
  \[
  0 \to H^{0}(X_{\mathrm{et}}, F') \to H^{0}(X_{\mathrm{et}}, F) \to H^{0}(X_{\mathrm{et}}, F'') \to H^{1}(X_{\mathrm{et}}, F') \to \cdots,
  \]
  而且长正合序列与短正合序列的关联是函子性的.
\end{enumerate}
此外，函子 $H^{r}(X_{\mathrm{et}}, \cdot)$ 由性质 (a),(b),(c) 唯一确定(到唯一的同构).

\begin{remark}\label{rem:9.1}
我们将经常使用如下论断：
设 $L = L_2 \circ L_1$，其中 $L_1$ 与 $L_2$ 都是带足够内射对象的 Abel 范畴之间的左正合函子.若 $L_1$ 保持内射对象且对某个对象 $X$，$(R^{r}L_1)(X) = 0$，则 $(R^{r}L)(X) = (R^{r}L_2)(L_1X)$.

证明是显然的：选取 $X$ 的内射分解 $X \to I^{\bullet}$，注意对 $L_1$ 的假设蕴含 $L_1X \to L_1I^{\bullet}$ 是 $L_1X$ 的内射分解，它可以用来计算 $(R^{r}L_2)(L_1X)$.现在 $(R^{r}L)(X)$ 与 $(R^{r}L_2)(L_1X)$ 都是 $L_2(L_1I^{\bullet}) = LI^{\bullet}$ 的第 $r$ 个上同调群.
\end{remark}

\begin{remark}\label{rem:9.2}
设 $\varphi \colon U \to X$ 是平展态射.如我们在注记~\ref{rem:8.16} 中所述，$\varphi^{*} \colon \mathrm{Sh}(X_{\mathrm{et}}) \to \mathrm{Sh}(U_{\mathrm{et}})$ 正合且保持内射对象.由于复合
\[
\mathrm{Sh}(X_{\mathrm{et}}) \xrightarrow{\varphi^{*}} \mathrm{Sh}(U_{\mathrm{et}}) \xrightarrow{\Gamma(U, \cdot)} \mathbf{Ab}
\]
是 $\Gamma(U, \cdot)$(回忆在这种情况下 $\varphi^{*}$ 就是限制)，我们看到 $F \mapsto F(U) \colon \mathrm{Sh}(X_{\mathrm{et}}) \to \mathbf{Ab}$ 的右导出函子是 $F \mapsto H^{r}(U_{\mathrm{et}}, F|_{U})$.我们常把 $H^{r}(U_{\mathrm{et}}, F|_{U})$ 写为 $H^{r}(U_{\mathrm{et}}, F)$.
\end{remark}

\begin{remark}\label{rem:9.3}
设 $\varphi \colon Y \to X$ 是态射.我们知道 $\varphi^{*}$ 正合(见注记~\ref{rem:8.9})，因此 $X$ 上层
\[
0 \to F' \to F \to F'' \to 0
\]
的短正合序列产生上同调群的长正合序列
\[
0 \to H^{0}(Y_{\mathrm{et}}, \varphi^{*}F') \to \cdots \to H^{r}(Y_{\mathrm{et}}, \varphi^{*}F') \to H^{r}(Y_{\mathrm{et}}, \varphi^{*}F) \to H^{r}(Y, \varphi^{*}F'') \to \cdots.
\]
由导出函子的泛性质(Weibel 2.4.7)，自然映射 $H^{0}(X_{\mathrm{et}}, F) \to H^{0}(Y_{\mathrm{et}}, \varphi^{*}F)$ 唯一延拓为与边缘映射相容的自然映射族 $H^{r}(X_{\mathrm{et}}, F) \to H^{r}(Y_{\mathrm{et}}, \varphi^{*}F)$.
\end{remark}

在本节余下部分，我们验证 Eilenberg-Steenrod 公理的类比成立.

\subsection*{维数公理}

设 $x = \operatorname{Spec} k$($k$ 是某个域), $\bar{x} = \operatorname{Spec} k^{\mathrm{sep}}$($k^{\mathrm{sep}}$ 是 $k$ 的某个可分闭包).如我们在 \S 6 中所述，函子
\[
F \mapsto M_F \overset{\mathrm{def}}{=} F_{\bar{x}}
\]
定义了 $x_{\mathrm{et}}$ 上层范畴与离散 $G$-模范畴之间的等价，其中 $G = \operatorname{Gal}(k^{\mathrm{sep}}/k)$.由于 $(M_F)^{G} = \Gamma(x, F)$，$M \mapsto M^{G}$ 与 $F \mapsto \Gamma(x, F)$ 的导出函子相对应.因此
\[
H^{r}(x, F) \cong H^{r}(G, M_F).
\]
为了使
\[
H^{r}(x, F) = 0 \quad (r > 0, \text{ 对所有 } F),
\]
正如维数公理所要求的，我们应该取 $x$ 为几何点，即可分闭域的谱.于是，正如我们在其他语境中已经看到的，对平展景而言，起拓扑空间中点的作用的是几何点.

\subsection*{正合性公理}

设 $Z$ 是 $X$ 的闭子簇(或子概形), $U = X \setminus Z$.对 $X_{\mathrm{et}}$ 上的任意层 $F$，定义
\[
\Gamma_Z(X, F) = \operatorname{Ker}(\Gamma(X, F) \to \Gamma(U, F)).
\]
于是 $\Gamma_Z(X, F)$ 是支集在 $Z$ 上的 $F$ 的截面群.我们有时从记号中去掉 $X$.函子 $F \mapsto \Gamma_Z(X, F)$ 显然左正合，我们把它对 $F$ 的第 $r$ 个右导出函子记为 $H^{r}_{Z}(X, \cdot)$($F$ 的支集在 $Z$ 上的上同调).

\begin{theorem}\label{thm:9.4}
对 $X_{\mathrm{et}}$ 上的任意层 $F$ 与闭子集 $Z \subset X$，有长正合序列
\[
\cdots \to H^{r}_{Z}(X, F) \to H^{r}(X, F) \to H^{r}(U, F) \to H^{r+1}_{Z}(X, F) \to \cdots.
\]
该序列关于空间对 $(X, X \setminus Z)$ 与 $F$ 函子性.
\end{theorem}

我们将在下一小节证明它.

\subsection*{Ext 群}

对固定的层 $F_0$，
\[
F \mapsto \operatorname{Hom}_X(F_0, F) \colon \mathrm{Sh}(X_{\mathrm{et}}) \to \mathbf{Ab}
\]
左正合，我们把它对 $F$ 的第 $r$ 个右导出函子记为 $\operatorname{Ext}^{r}(F_0, \cdot)$.于是，
\begin{enumerate}
  \item 对任意层 $F$，$\operatorname{Ext}^{0}(F_0, F) = \operatorname{Hom}(F_0, F)$；
  \item 若 $\mathcal{I}$ 内射，则对 $r > 0$，$\operatorname{Ext}^{r}(F_0, \mathcal{I}) = 0$；
  \item 层的短正合序列
  \[
  0 \to F' \to F \to F'' \to 0
  \]
  产生长正合序列
  \[
  \cdots \to \operatorname{Ext}^{r}_{X}(F_0, F') \to \operatorname{Ext}^{r}_{X}(F_0, F) \to \operatorname{Ext}^{r}_{X}(F_0, F'') \to \cdots,
  \]
  而且长正合序列与短正合序列的关联是函子性的.
\end{enumerate}

\begin{example}\label{ex:9.5}
设 $\mathbb{Z}$ 表示 $X$ 上的常值层.对 $X$ 上的任意层 $F$，映射 $\alpha \mapsto \alpha(1)$ 是同构 $\operatorname{Hom}_X(\mathbb{Z}, F) \xrightarrow{\sim} F(X)$.因此 $\operatorname{Hom}_X(\mathbb{Z}, \cdot) \cong \Gamma(X, \cdot)$，从而 $\operatorname{Ext}^{r}_{X}(\mathbb{Z}, \cdot) \cong H^{r}(X_{\mathrm{et}}, \cdot)$.

由于 $\operatorname{Hom}_X(F_0, \cdot)$ 关于 $F_0$ 函子性，$\operatorname{Ext}^{r}_{X}(F_0, \cdot)$ 也关于 $F_0$ 函子性.
\end{example}

\begin{proposition}\label{prop:9.6}
$X_{\mathrm{et}}$ 上层
\[
0 \to F''_0 \to F_0 \to F'_0 \to 0
\]
的短正合序列对任意层 $F$ 产生长正合序列
\[
\cdots \to \operatorname{Ext}^{r}_{X}(F'_0, F) \to \operatorname{Ext}^{r}_{X}(F_0, F) \to \operatorname{Ext}^{r}_{X}(F''_0, F) \to \cdots.
\]
\end{proposition}

\begin{proof}
若 $\mathcal{I}$ 内射，则
\[
0 \to \operatorname{Hom}_X(F''_0, \mathcal{I}) \to \operatorname{Hom}_X(F_0, \mathcal{I}) \to \operatorname{Hom}_X(F'_0, \mathcal{I}) \to 0
\]
正合.对 $F$ 的任何内射分解 $F \to \mathcal{I}^{\bullet}$，
\[
0 \to \operatorname{Hom}_X(F''_0, \mathcal{I}^{\bullet}) \to \operatorname{Hom}_X(F_0, \mathcal{I}^{\bullet}) \to \operatorname{Hom}_X(F'_0, \mathcal{I}^{\bullet}) \to 0
\]
是复形的正合序列，由同调代数的标准结果，它产生上同调群的长正合序列
\[
\cdots \to \operatorname{Ext}^{r}_{X}(F''_0, F) \to \operatorname{Ext}^{r}_{X}(F_0, F) \to \operatorname{Ext}^{r}_{X}(F'_0, F) \to \cdots.
\]
\end{proof}

现在证明定理~\ref{thm:9.4}.设
\[
U \xrightarrow{j} X \xleftarrow{i} Z
\]
如定理的叙述中所示.设 $\mathbb{Z}$ 表示由 $\mathbb{Z}$ 定义的 $X$ 上的常值层，考虑正合序列 (\ref{prop:8.15})
\[
0 \to j_{!}j^{*}\mathbb{Z} \to \mathbb{Z} \to i_{*}i^{*}\mathbb{Z} \to 0. \qquad (*)
\]
对 $X_{\mathrm{et}}$ 上的任意层 $F$，
\[
\operatorname{Hom}_X(j_{!}j^{*}\mathbb{Z}, F) = \operatorname{Hom}_U(j^{*}\mathbb{Z}, j^{*}F) = F(U),
\]
从而 $\operatorname{Ext}^{r}_{X}(j_{!}j^{*}\mathbb{Z}, F) = H^{r}(U_{\mathrm{et}}, F)$.由正合序列
\[
0 \to \operatorname{Hom}(i_{*}i^{*}\mathbb{Z}, F) \to \operatorname{Hom}(\mathbb{Z}, F) \to \operatorname{Hom}(j_{!}j^{*}\mathbb{Z}, F)
\]
发现 $\operatorname{Hom}_X(i_{*}i^{*}\mathbb{Z}, F) = \Gamma_Z(X, F)$，从而 $\operatorname{Ext}^{r}_{X}(i_{*}i^{*}\mathbb{Z}, F) = H^{r}_{Z}(X, F)$.因此，对应于 $(*)$，如 (\ref{prop:9.6}) 中的 Ext 的长正合序列，就是定理~\ref{thm:9.4} 所需的序列.

\subsection*{切除}

拓扑空间的切除说，支集在 $Z$ 上的上同调应当只依赖于 $Z$ 在 $X$ 中的邻域，例如，把 $X$ 换成 $Z$ 的开邻域不应改变 $H^{r}_{Z}(X, F)$.下面是对平展拓扑的类似陈述.

\begin{theorem}\label{thm:9.7}[切除]
设 $\varphi \colon X' \to X$ 是平展映射，$Z' \subset X'$ 是 $X'$ 的闭子簇(或子概形)，使得
\begin{enumerate}
  \item $Z \overset{\mathrm{def}}{=} \varphi(Z')$ 在 $X$ 中闭，且 $\varphi$ 到 $Z'$ 的限制是同构 $Z' \to Z$；
  \item $\varphi(X' \setminus Z') \subset X \setminus Z$.
\end{enumerate}
则对 $X_{\mathrm{et}}$ 上的任意层 $F$，典范映射 $H^{r}_{Z}(X_{\mathrm{et}}, F) \to H^{r}_{Z'}(X'_{\mathrm{et}}, F|_{X'})$ 对所有 $r$ 都是同构.
\end{theorem}

\begin{proof}
设 $U' = X' \setminus Z'$，$U = X \setminus Z$.我们有交换图
\[
\begin{tikzcd}[row sep=small, column sep=small]
  U' \arrow[r, "j'"] \arrow[d, "\varphi"] & X' \arrow[d, "\varphi"] & Z' \arrow[l, "i'"'] \arrow[d, "\varphi"] \\
  U \arrow[r, "j"'] & X & Z \arrow[l, "i"]
\end{tikzcd}
\]
这产生图
\[
\begin{tikzcd}
0 \arrow[r] & \Gamma_{Z'}(X', \varphi^{*}F) \arrow[r] & \Gamma(X', \varphi^{*}F) \arrow[r] & \Gamma(U', \varphi^{*}F) \\
0 \arrow[r] & \Gamma_{Z}(X, F) \arrow[r] \arrow[u] & \Gamma(X, F) \arrow[r] \arrow[u] & \Gamma(U, F) \arrow[u]
\end{tikzcd}
\]
第一个竖箭头由其余两个诱导.由于 $\varphi^{*}$ 正合且保持内射对象，只需对 $r = 0$ 证明定理，即证明上图 (对 $F$) 中的映射 $\Gamma_Z(X, F) \to \Gamma_{Z'}(X', F|_{X'})$ 是同构.

设 $s \in \Gamma_Z(X, F)$ 在 $\Gamma_{Z'}(X', F|_{X'})$ 中映到零.则 $s$(看作 $\Gamma(X, F)$ 的元素)在 $\Gamma(X', F)$ 与 $\Gamma(U, F)$ 中限制为零.由于映射对 $(X' \to X, U \to X)$ 是 $X$ 的覆盖且 $F$ 是层，这蕴含 $s = 0$.

设 $s' \in \Gamma_{Z'}(X', F|_{X'})$，把它看作 $\Gamma(X', F)$ 的元素.我们必须找到限制为 $s'$(在 $\Gamma(X', F)$ 中)与 $0$(在 $\Gamma(U, F)$ 中)的 $s \in \Gamma(X, F)$.一旦我们检验了两个截面在``重叠''处一致，层判别准则就为我们提供这样的截面.首先注意 $U \times_{X} X' = \varphi^{-1}(U) = U'$，且两个截面在 $U'$ 上都限制为零.还剩下检验 $s'$ 在两个映射 $X' \rightrightarrows X' \times_{X} X'$ 下的限制相等，这可以在茎上做.对 $U' \times_{X} U'$ 中的点，两个限制都为零.两个映射 $Z' \rightrightarrows Z' \times_{X} Z'$ 相等，因此在 $Z' \times_{X} Z'$ 的点处两个限制一致.由于 $X' \times_{X} X' = U' \times_{X} U' \cup Z' \times_{X} Z'$，这完成了证明.
\end{proof}

定理中的情形出现于：$\varphi \colon X' \to X$ 平展，且 $X$ 含闭子概形，使得存在态射 $s \colon Z \to X'$，满足 $\varphi \circ s = \mathrm{id}_{Z}$ 且 $\varphi^{-1}(Z) = s(Z)$.

\begin{corollary}\label{cor:9.8}
设 $x$ 是 $X$ 的闭点.对 $X$ 上的任意层 $F$，有同构 $H^{r}_{x}(X, F) \to H^{r}_{x}(\operatorname{Spec} \mathcal{O}^{\mathrm{h}}_{X, x}, F)$，其中 $\mathcal{O}^{\mathrm{h}}_{X, x}$ 是 $\mathcal{O}_{X, x}$ 的 Hensel 化.
\end{corollary}

\begin{proof}
按定理，对 $x$ 的任何使 $u$ 是 $U$ 中唯一映到 $x$ 的点的平展邻域 $(U, u)$，$H^{r}_{x}(X, F) = H^{r}_{u}(U, F)$.这样的平展邻域是共尾的，对它们取极限，得到所需的同构(关于取概形的逆向极限时上同调的行为见下面 10.9).
\end{proof}

\subsection*{同伦公理}

回忆同伦公理说同伦的映射在上同调上诱导相同的映射.这个陈述有一个类比，其中同伦等价换成有理等价.

为简单起见，设 $X$ 是代数闭域 $k$ 上的代数簇.设 $Z$ 是 $X \times \mathbb{P}^{1}$ 的闭子簇，它在投影映射 $\pi \colon X \times \mathbb{P}^{1} \to \mathbb{P}^{1}$ 下的像在 $\mathbb{P}^{1}$ 中稠密.对任意闭点 $t \in \mathbb{P}^{1}$，$Z(t) \overset{\mathrm{def}}{=} Z \cap (X \times \{t\})$ 是 $X \times \{t\} = X$ 的闭子簇(相交应在相交理论的意义下进行，即允许重数---见后面).我们可以把 $t \in \pi(Z)$ 的 $Z(t)$ 看作 $X$ 上的代数循环的连续族，并且如果 $Z_1$ 与 $Z_2$ 出现在这样的族中，我们写 $Z_1 \sim Z_2$.我们说两个代数循环 $Z$ 与 $Z'$ 有理等价，如果存在代数循环 $Z_1, Z_2, \ldots, Z_r$，使得
\[
Z \sim Z_1 \sim \cdots \sim Z_r \sim Z'.
\]

\begin{theorem}\label{thm:9.9}
若两个态射 $\varphi, \varphi'$ 的图象有理等价，则它们在平展上同调上定义相同的映射.
\end{theorem}

\begin{proof}
态射定义的上同调映射只依赖于态射图象 $\Gamma_{\varphi}$ 的上同调类，而这个上同调类只依赖于 $\Gamma_{\varphi}$ 的有理等价类(见后面 \S 23).
\end{proof}
